\documentclass{article}
\usepackage[utf8]{inputenc}
\usepackage[french]{babel}
\usepackage{geometry}
\usepackage{listings}
\usepackage{xcolor}
\usepackage{hyperref}
\usepackage{graphicx}

\geometry{a4paper, margin=2cm}

\definecolor{codegreen}{rgb}{0,0.6,0}
\definecolor{codegray}{rgb}{0.5,0.5,0.5}
\definecolor{codepurple}{rgb}{0.58,0,0.82}
\definecolor{backcolour}{rgb}{0.95,0.95,0.92}

\lstdefinestyle{mystyle}{
    backgroundcolor=\color{backcolour},
    commentstyle=\color{codegreen},
    keywordstyle=\color{magenta},
    numberstyle=\tiny\color{codegray},
    stringstyle=\color{codepurple},
    basicstyle=\ttfamily\footnotesize,
    breakatwhitespace=false,
    breaklines=true,
    captionpos=b,
    keepspaces=true,
    numbers=left,
    numbersep=5pt,
    showspaces=false,
    showstringspaces=false,
    showtabs=false,
    tabsize=2
}

\lstset{style=mystyle}

\title{Analyse Complète du Backend Unit Converter}
\author{Votre Assistant IA}
\date{\today}

\begin{document}

\maketitle

\tableofcontents

\section{Introduction et Structure Globale}

Ce document détaille l'architecture backend de votre projet \texttt{Unit Converter}, construit avec \textbf{Java Spring Boot}. Ce projet suit une architecture en couches (Layered Architecture) classique et robuste, idéale pour les applications d'entreprise et les projets maintenables.

\subsection{Objectif du Backend}

Le backend sert de "cerveau" à votre application :
\begin{itemize}
    \item Il expose une \textbf{API REST} (Representational State Transfer) que le frontend consomme.
    \item Il gère la \textbf{logique métier} complexe (les conversions d'unités).
    \item Il assure la \textbf{sécurité} (authentification, tokens JWT, rôle des utilisateurs).
    \item Il persiste les données dans une base de données \textbf{MySQL} via \textbf{Hibernate/JPA}.
\end{itemize}

\subsection{Structure des Packages}

Votre code source principal se trouve dans \texttt{src/main/java/fullstack/unit\_converter\_backend}. Voici l'organisation des packages et leur rôle :

\begin{description}
    \item[controller] Contient les classes qui reçoivent les requêtes HTTP (les points d'entrée de l'API). C'est la couche de présentation.
    \item[service] Contient la logique métier pure (les calculs, les règles de gestion). C'est le cœur de l'application.
    \item[repository] Contient les interfaces pour interagir avec la base de données. C'est la couche d'accès aux données (DAO - Data Access Object).
    \item[model] (ou entity) Contient les classes Java qui représentent les tables de votre base de données.
    \item[dto] (Data Transfer Object) Contient des objets simples utilisés pour transporter des données entre le client (frontend) et le serveur, sans exposer directement les modèles de base de données.
    \item[security] Contient toute la configuration liée à la sécurité (JWT, filtres, configuration Spring Security).
    \item[enums] Contient les énumérations (constantes typées) pour les unités, les rôles, etc.
\end{description}

\section{Point d'Entrée de l'Application}

\subsection{UnitConverterBackendApplication.java}

C'est le fichier qui lance tout le processus.

\begin{lstlisting}[language=Java]
package fullstack.unit_converter_backend;

import org.springframework.boot.SpringApplication;
import org.springframework.boot.autoconfigure.SpringBootApplication;

@SpringBootApplication
public class UnitConverterBackendApplication {

    public static void main(String[] args) {
        SpringApplication.run(UnitConverterBackendApplication.class, args);
    }
}
\end{lstlisting}

\subsubsection{Explication Détailée}

\begin{itemize}
    \item \textbf{@SpringBootApplication} : C'est une "méta-annotation" qui combine trois annotations importantes :
    \begin{enumerate}
        \item \texttt{@Configuration} : Indique que cette classe peut définir des beans (objets gérés par Spring).
        \item \texttt{@EnableAutoConfiguration} : Dit à Spring Boot de configurer automatiquement l'application en fonction des dépendances présentes (ex: s'il voit \texttt{spring-boot-starter-web}, il lance un serveur Tomcat; s'il voit \texttt{mysql-connector}, il configure la connexion BDD).
        \item \texttt{@ComponentScan} : Dit à Spring de scanner le package actuel (\texttt{fullstack.unit\_converter\_backend}) et tous ses sous-packages pour trouver vos \texttt{@Controller}, \texttt{@Service}, \texttt{@Repository}, etc. \textbf{C'est pourquoi la position de ce fichier à la racine du code source est cruciale.}
    \end{enumerate}

    \item \textbf{public static void main(String[] args)} : Le point d'entrée standard de toute application Java.

    \item \textbf{SpringApplication.run(...)} : Cette méthode statique démarre le conteneur Spring (ApplicationContext). Elle lance le serveur web embarqué (Tomcat) sur le port défini (8080 par défaut).
\end{itemize}

\section{Architecture et Flux d'une Requête}

Pour comprendre comment le backend fonctionne, suivons le cheminement d'une requête HTTP typique (ex: \texttt{POST /api/conv/preview}).

\begin{enumerate}
    \item \textbf{Client (Frontend)} : Envoie une requête HTTP avec des données JSON.
    \item \textbf{Security Filter Chain} : La requête est interceptée par Spring Security (\texttt{JwtAuthFilter}). On vérifie si l'utilisateur est authentifié (token valide).
    \item \textbf{Controller} (\texttt{ConversionController}) : Reçoit la requête, valide les données (ex: champs non vides), et appelle le service.
    \item \textbf{Service} (\texttt{ConversionService}) : Exécute la logique métier (le calcul de conversion).
    \item \textbf{Repository} (si besoin) : Si on devait sauvegarder l'historique, le service appellerait le repository.
    \item \textbf{Base de Données} : Les données sont lues ou écrites.
    \item \textbf{Retour} : Le résultat remonte la chaîne : BDD $\rightarrow$ Repository $\rightarrow$ Service $\rightarrow$ Controller.
    \item \textbf{Controller} : Renvoie un objet \texttt{ResponseEntity} contenant la réponse (DTO) et le code HTTP (ex: 200 OK).
    \item \textbf{Jackson} : Spring transforme automatiquement l'objet Java en JSON pour le client.
\end{enumerate}

\section{Analyse des Contrôleurs (Controllers)}

Les contrôleurs sont les "réceptionnistes" de votre application.

\subsection{Annotations Communes}
\begin{itemize}
    \item \textbf{@RestController} : Indique que cette classe gère des requêtes REST et que chaque méthode renvoie directement la réponse (souvent du JSON) dans le corps de la réponse HTTP (pas de vue HTML).
    \item \textbf{@RequestMapping("/api/...")} : Définit le préfixe d'URL commun à toutes les méthodes de la classe.
    \item \textbf{@GetMapping, @PostMapping} : Mappe une méthode Java à une méthode HTTP spécifique (GET pour lire, POST pour créer/agir).
    \item \textbf{@RequestBody} : Dit à Spring de prendre le JSON envoyé par le client et de le transformer en objet Java (désérialisation).
    \item \textbf{@Valid} : Demande à Spring de vérifier les contraintes de validation (ex: \texttt{@NotBlank}) définies dans le DTO avant d'exécuter la méthode.
\end{itemize}

\subsection{AuthController.java}
Ce contrôleur gère l'inscription et la connexion.

\begin{itemize}
    \item \textbf{Endpoints} :
    \begin{itemize}
        \item \texttt{POST /api/auth/register} : Crée un nouvel utilisateur. Appelle \texttt{userService.register()}.
        \item \texttt{POST /api/auth/login} : Authentifie un utilisateur. Utilise \texttt{authenticationManager.authenticate()} pour vérifier l'email/mot de passe, puis génère un token JWT.
    \end{itemize}
\end{itemize}

\subsection{ConversionController.java}
Gère les conversions d'unités.

\begin{itemize}
    \item \textbf{Endpoints} :
    \begin{itemize}
        \item \texttt{GET /api/conv/info} : Simple test de disponibilité.
        \item \texttt{POST /api/conv/preview} : Calcule une conversion sans la sauvegarder (ou la sauvegarde, selon votre logique actuelle).
        \\ \textbf{Logique importante} : Elle reçoit un \texttt{ConversionRequest}, appelle le service de conversion, récupère l'utilisateur connecté via \texttt{Authentication}, et sauvegarde l'historique via \texttt{HistoryService}.
        \item \texttt{GET /api/conv/history} : Récupère l'historique de l'utilisateur connecté.
    \end{itemize}
\end{itemize}

\subsection{UserController.java}
Permet de gérer et consulter les utilisateurs (souvent pour un tableau de bord administrateur).

\begin{itemize}
    \item \textbf{Endpoints} :
    \begin{itemize}
        \item \texttt{GET /api/users} : Renvoie la liste de tous les utilisateurs (via \texttt{UserResponse}).
        \item \texttt{GET /api/users/{id}/history} : Renvoie l'historique d'un utilisateur specifique.
        \item \textbf{Note} : Ces routes sont généralement protégées et accessibles uniquement aux administrateurs (rôle ADMIN), configuré dans \texttt{SecurityConfig} (\texttt{requestMatchers("/api/users/**").hasRole("ADMIN")}).
    \end{itemize}
\end{itemize}

\section{Analyse des Services}

Les services contiennent le "comment" de l'application.

\subsection{UserService.java}
Gère les utilisateurs et l'authentification.
\begin{itemize}
    \item Implémente \textbf{UserDetailsService} : C'est une interface standard de Spring Security. La méthode \texttt{loadUserByUsername} est cruciale : elle charge un utilisateur depuis la BDD pour que Spring Security puisse vérifier le mot de passe.
    \item \textbf{register()} : Vérifie si l'email existe, chiffre le mot de passe (\texttt{passwordEncoder}), sauvegarde l'utilisateur, et génère un token JWT initial.
    \item \textbf{Dépendances} : \texttt{UserRepository} (accès BDD), \texttt{PasswordEncoder} (sécurité), \texttt{JwtUtil} (tokens).
\end{itemize}

\subsection{ConversionService.java}
Contient la logique pure de conversion.
\begin{itemize}
    \item Pas de dépendance vers la base de données ici, c'est du pur calcul.
    \item Utilise des méthodes privées (\texttt{convertWeight}, \texttt{convertDistancce}...) pour organiser le code.
    \item Lance des exceptions (\texttt{IllegalArgumentException}) si les unités sont invalides.
\end{itemize}

\subsection{HistoryService.java}
Gère la persistance et la gestion des historiques de conversion.
\begin{itemize}
    \item \textbf{saveHistory} : Prend le résultat d'une conversion, crée un objet \texttt{ConversionHistory}, et le sauvegarde avec \texttt{repo.save()}.
    \item \textbf{getUserHistory} : Récupère les conversions d'un utilisateur, en utilisant un repository avec une méthode qui trie par date (\texttt{OrderByCreatedAtDesc}).
    \item \textbf{Transformation DTO} : La méthode privée \texttt{toDto()} assure la conversion entre l'entité BDD et le \texttt{HistoryResponse} renvoyé à l'API. C'est essentiel pour ne pas exposer l'entité brute.
\end{itemize}

\section{Modèles et Données (Layer Entity & Repository)}

\subsection{Les Entités (Entities)}
Ce sont les reflets Java de vos tables SQL.

\textbf{User.java} :
\begin{itemize}
    \item \texttt{@Entity}, \texttt{@Table(name="users")} : Mappe la classe à la table `users`.
    \item \texttt{@Id}, \texttt{@GeneratedValue} : Clé primaire auto-incrémentée.
    \item \texttt{@OneToMany(mappedBy="user")} : Relation inverse. Un utilisateur a plusieurs historiques. \texttt{mappedBy} indique que la clé étrangère est dans l'autre table (`conversion_history`).
\end{itemize}

\textbf{ConversionHistory.java} :
\begin{itemize}
    \item \texttt{@ManyToOne} : Relation directe. Un historique appartient à un seul utilisateur.
    \item \texttt{@JoinColumn(name="user\_id")} : Définit la colonne de clé étrangère dans la table `conversion\_history`.
\end{itemize}

\subsection{Les Repositories}
Interfaces qui étendent \texttt{JpaRepository}.
\begin{itemize}
    \item Spring Data JPA génère automatiquement l'implémentation de ces interfaces au démarrage.
    \item Vous n'avez pas besoin d'écrire de SQL pour des opérations simples (save, findById, delete).
    \item \textbf{Méthodes dérivées} : \texttt{findByUserOrderByCreatedAtDesc(User user)}
    \\ Spring comprend le nom de la méthode et génère une requête SQL du type :
    \\ \texttt{SELECT * FROM conversion\_history WHERE user\_id = ? ORDER BY created\_at DESC}
\end{itemize}

\section{Sécurité (Spring Security & JWT)}

Votre sécurité est "Stateless" (sans session serveur), basée sur des tokens JWT.

\subsection{Le Flux d'Authentification}
\begin{enumerate}
    \item L'utilisateur envoie ses identifiants via \texttt{/api/auth/login}.
    \item Si valide, le serveur renvoie un \textbf{Token JWT} (chaîne cryptée avec une clé secrète).
    \item Pour les requêtes suivantes, le client doit envoyer ce token dans le header HTTP : \texttt{Authorization: Bearer <token>}.
\end{enumerate}

\subsection{JwtAuthFilter.java}
Ce filtre s'exécute avant chaque requête (sauf celles publiques).
\begin{itemize}
    \item Il intercepte la requête.
    \item Il regarde le header \texttt{Authorization}.
    \item Il extrait le token et le valide via \texttt{JwtUtil}.
    \item Si valide, il crée un objet \texttt{Authentication} et le place dans le \texttt{SecurityContext}. Ainsi, Spring sait qui est l'utilisateur pour le reste de la requête (permettant d'utiliser \texttt{Authentication authentication} dans vos contrôleurs).
\end{itemize}

\section{Configuration et Base de Données}

\subsection{application.properties}
\begin{lstlisting}
spring.datasource.url=jdbc:mysql://localhost:3306/unit_converter_db...
spring.jpa.hibernate.ddl-auto=update
\end{lstlisting}
\begin{itemize}
    \item \textbf{ddl-auto=update} : Très utile en dev. Au démarrage, Hibernate regarde vos classes \texttt{@Entity} et met à jour automatiquement le schéma de la base de données (crée les tables, ajoute les colonnes manquantes) sans supprimer les données existantes.
\end{itemize}

\section{DTOs (Data Transfer Objects)}

\subsection{Pourquoi des DTOs ?}
Pourquoi ne pas utiliser directement l'entité \texttt{User} dans le contrôleur ?
\begin{itemize}
    \item \textbf{Sécurité} : L'entité \texttt{User} contient le mot de passe hashé. On ne veut jamais renvoyer ça au client. Le \texttt{UserResponse} ne contient que les infos publiques.
    \item \textbf{Découplage} : La structure de votre base de données (Entity) peut changer sans casser l'API (DTO).
    \item \textbf{Validation} : Les annotations de validation (\texttt{@NotBlank}, \texttt{@Email}) sont placées sur les DTOs d'entrée (\texttt{RegistrationRequest}) pour valider les données avant même qu'elles n'atteignent le service.
\end{itemize}

\section{Conclusion}

Votre backend est bien structuré. Il sépare clairement les responsabilités, utilise les standards actuels (Spring Boot 3+, Spring Security avec JWT, JPA), et suit les bonnes pratiques de développement (DTOs, Services sans état, Injection de dépendances par constructeur).

\end{document}



